\begin{table}[tbp]
\centering
\caption{\textbf{How you look decides what turns up.} Three ways of letting last
year's finding change this year's audit, with the underlying problems held
identical and drawn independently each period. Each cell is the change in the
gap between programs with and without a prior finding, measured against a
baseline in which prior findings change nothing. Following up the same finding
moves repeats and \emph{only} repeats. Looking harder across a requirement area
also surfaces new findings inside that area. Looking harder at the program
overall also surfaces findings in areas that never had one --- which is what an
auditor would otherwise read as a genuinely new problem. The zeros are exact:
they come from using the same random draws across regimes.}
\label{tab:multitype}
\small
\begin{tabular}{lccc}
\toprule
& \multicolumn{3}{c}{Change in the prior-finding gap} \\
\cmidrule(lr){2-4}
If last year's finding makes the auditor\dots
& \shortstack{A finding\\repeats}
& \shortstack{Something new,\\in an area that\\already had one}
& \shortstack{Something new,\\in an area that\\did not} \\
\midrule
Follow up the same finding & $+0.0394$ & $+0.0000$ & $+0.0000$ \\
Look harder at the same requirement area & $+0.0279$ & $+0.0454$ & $+0.0000$ \\
Look harder at the program overall & $+0.0223$ & $+0.0366$ & $+0.1388$ \\
\bottomrule
\end{tabular}
\end{table}
